% !TeX root = ../../negocio-tecnologico.tex
% Modelo para preenchimento: substitua as orientacoes por seu relato final.
\chapter{Introdução e oportunidade}
Apresente problema, demanda, público, contexto, justificativa, objetivos e
recorte do negócio estudado. Identifique a base tecnológica. O roteiro
organiza uma entrega final de TCC, conforme o manual \citeind{manualundf},
sem pressupor abertura de empresa ou faturamento.

\chapter{Fundamentação e método}
Articule as referências e descreva os procedimentos de investigação e análise
efetivamente empregados. Informe fontes, instrumentos, critérios e limites.
Ferramentas de modelagem devem ser justificadas, não tratadas como exigências
universais. Cadastre as fontes no arquivo bibliográfico da entrada principal.

\chapter{Análise do contexto e do mercado}
Apresente evidências sobre público, necessidades, alternativas e condições
relevantes à proposta. Identifique a origem e o alcance dos dados. Se não
realizou pesquisa com clientes, declare essa limitação e não apresente
hipóteses como aceitação de mercado comprovada.

\chapter{Solução e base tecnológica}
Descreva produto ou serviço, tecnologia, diferenciação e estágio alcançado.
Distinga conceito, protótipo, piloto e operação. Identifique os artefatos
produzidos e as dependências ainda não resolvidas.

\chapter{Modelo de negócio consolidado}
Explique proposta de valor, segmentos, canais, relacionamento, atividades,
recursos, parcerias, receitas e custos e suas relações. Quadros e ferramentas
visuais são sínteses que precisam de discussão. Registre quais componentes
têm evidência e quais continuam como hipóteses.

\chapter{Resultados da investigação e avaliação}
Relate achados, testes ou análises realizados e como alteraram a proposta.
Identifique hipóteses examinadas e não examinadas. Não invente entrevistas,
clientes, vendas ou métricas para preencher o modelo. Apresente análises
negativas ou inconclusivas com suas evidências e limites.

\chapter{Viabilidade, riscos e limitações}
Analise a viabilidade técnica, operacional e econômico-financeira dentro do
escopo. Explicite fontes, premissas, cálculos, cenários e incertezas. Identifique
projeções de demanda, custos e receita como estimativas. Não confunda equilíbrio
aritmético de um cenário com viabilidade comercial comprovada.

\chapter{Conclusão}
Responda aos objetivos com posicionamento fundamentado sobre a proposta.
Sintetize contribuição, limitações, condições de aplicação e questões abertas.
Recomendações futuras devem decorrer da análise do trabalho concluído, sem
prometer sucesso empresarial ou substituir o relatório por um cronograma.
